\VerdierPhysicalPageAnchor{0022}%
up to homotopy by the complex $Y^{\bullet}$. The cohomology objects
$\mathrm R^nF(Y^{\bullet})$ of the complex
$\int F(I^{\bullet\bullet})$ are therefore defined up to canonical
isomorphism, and one sees that they depend functorially on the complex
$Y^{\bullet}$. The resulting functors $\mathrm R^nF$ are the
\emph{right hyperderived functors} of $F$.

\VerdierSemanticAnchor{intro-2-spectral-sequences}
One also obtains two spectral sequences, functorial in $Y^{\bullet}$, whose
initial terms are:

\VerdierSemanticAnchor{intro-2-spectral-sequence-one}
\begin{equation}
\mathrm I_2^{p,q}=\mathrm H^p\bigl(\mathrm R^qF(Y^{\bullet})\bigr)
\tag{*}
\end{equation}

\VerdierSemanticAnchor{intro-2-spectral-sequence-two}
\begin{equation}
\mathrm{II}_2^{p,q}=\mathrm R^pF\bigl(\mathrm H^q(Y^{\bullet})\bigr),
\tag{**}
\end{equation}
and whose limiting terms are the graded objects associated with the suitably
filtered objects $\mathrm R^nF(Y^{\bullet})$. For simplicity, let us restrict
to the case in which the components of the complexes under consideration
vanish in sufficiently small degree. One then sees that the spectral
sequences (*) and (**) converge. A morphism
$m:Y^{\bullet}\to Z^{\bullet}$ of complexes induces morphisms

\VerdierSemanticAnchor{intro-2-hyperderived-morphism}
\[
\mathrm R^nF(m):\mathrm R^nF(Y^{\bullet})
\longrightarrow \mathrm R^nF(Z^{\bullet})
\]
and morphisms of the corresponding spectral sequences.

\VerdierSemanticAnchor{intro-2-quasi-isomorphism-invariance}
In particular, \emph{when $m$ is a quasi-isomorphism}, spectral sequence (**)
shows that \emph{the morphisms $\mathrm R^nF(m)$ are isomorphisms}. Thus the
objects $\mathrm R^nF(Y^{\bullet})$ are determined once the complex
$Y^{\bullet}$ is known ``up to quasi-isomorphism''. Since, moreover, we have
observed that the usual methods of homological algebra most often furnish
only complexes determined up to quasi-isomorphism, it appears that the
natural objects studied by hyperhomology are the ``classes of complexes up to
quasi-isomorphism''.

\VerdierSemanticAnchor{intro-3}
\noindent\textbf{3.}\quad
The approach we follow is thus dictated naturally. For every abelian category
$\mathcal A$, we introduce the category $\operatorname{comp}(\mathcal A)$ of
complexes of objects of $\mathcal A$. We then construct from
$\operatorname{comp}(\mathcal A)$ a new category by formally inverting the
quasi-isomorphisms of $\operatorname{comp}(\mathcal A)$.

\VerdierSemanticAnchor{intro-3-derived-category}
We thereby obtain a category $\mathbf D(\mathcal A)$: the
\emph{``derived category'' of the abelian category $\mathcal A$}. By means of
right resolutions, which need not be specified in this introduction, one
associates to an additive functor between two abelian categories
$F:\mathcal A\to\mathcal B$ a functor
$\mathrm RF:\mathbf D(\mathcal A)\to\mathbf D(\mathcal B)$ (which most often
is defined only on a suitable subcategory
$\mathbf D^*(\mathcal A)\subset\mathbf D(\mathcal A)$, obtained by imposing
bounds on the degrees of the complexes considered), called the
\emph{total right derived functor} of $F$.

\VerdierSemanticAnchor{intro-3-total-right-derived-functor}
For every object $X$ of $\mathbf D^*(\mathcal A)$, the object $\mathrm RF(X)$
of $\mathbf D(\mathcal B)$ is thus a complex of objects of $\mathcal B$
defined ``up to quasi-isomorphism''; consequently, the cohomology objects
$\mathrm R^nF(X)$ of this complex are well defined. The functors
$X\longmapsto\mathrm R^nF(X)$ are the ``right hyperderived functors'' of $F$,
defined here without necessarily assuming that
